\documentclass[11pt]{article}
\usepackage{../../../tex/lecturestyle}

\title{Matchings on a Path and Their Applications}
\author{}
\date{}

\begin{document}
\thispagestyle{fancy}

\lecturetitleband{Graph Theory}{Matchings on a Path and Their Applications}{8}

\section{Overview}
Three problems that look unrelated turn out to be the same problem.

\begin{itemize}[leftmargin=1.5em,itemsep=1pt]
  \item \emph{Non-adjacent picks.} Choose exactly $K$ entries of an array,
        no two of them neighbours, with the largest possible sum.
  \item \emph{Disjoint subarrays.} Choose exactly $K$ pairwise disjoint
        non-empty subarrays of an array with the largest possible total.
  \item \emph{Stock trading.} Given the prices on days $0,\dots,N$, make
        exactly $K$ trades (buy, then sell later, holding one share at a
        time) with the largest possible profit.
\end{itemize}

Each has a well-known greedy solution that raises $K$ one step at a time
and repairs the previous optimum locally. Such greedy methods are usually
presented as tricks to be memorised. Here they are \emph{derived}: all three
problems become a maximum-weight matching problem on a path
(Sections~\ref{sec:path} and~\ref{sec:reductions}), and a general theorem
about augmenting paths (Section~\ref{sec:greedy}) produces the greedy
algorithm, with its correctness proof, in one stroke. A closer look at
the order in which the greedy takes its decisions
(Section~\ref{sec:linear}) gives a stack-based $O(n)$ algorithm for one
fixed $K$. The last part (Section~\ref{sec:range}) goes further and
answers \emph{static range queries} --- the best size-$K$ matching on a
sub-interval --- in $O(\log n)$ each.

\paragraph{Prerequisites.} Basic notions about matchings (covered vertex,
alternating path) and a priority queue. Section~\ref{sec:range} also uses
a Fenwick tree, a range-minimum structure, a segment tree and a wavelet
matrix (see the Wavelet Matrix lectures in this series for the latter).

\paragraph{Summary of results.} Let $n$ be the number of vertices.
\begin{center}\small\sffamily
\begin{tabular}{@{}ll@{}}
\toprule
task & time\\
\midrule
all values $f(0),f(1),\dots$ (heap greedy) & $O(n\log n)$\\
all values, via the stack and one sort & $O(n\log n)$, small constant\\
a single value $f(K)$ (stack + selection) & $O(n)$\\
$q$ range queries $(L,R,K)$, offline & $O((n+q)\log n)$\\
\bottomrule
\end{tabular}
\end{center}

\section{Matchings on a Path}\label{sec:path}
Throughout, the \emph{path} on $n$ vertices has vertices $0,1,\dots,n-1$ and
edges $\{i,i+1\}$. A \emph{matching} is a set of pairwise disjoint edges;
a vertex on one of its edges is \emph{covered}. A matching of size $k$
covers $2k$ vertices, and a matching of size $k$ exists iff
$0\le k\le\lfloor n/2\rfloor$.

\begin{definitionbox}[Vertex-weighted path matching]
Every vertex $i$ carries a weight $x_i\in\mathbb Z$. The \emph{value} of a
matching is the total weight of the vertices it covers, and
\[
  f_x(k)\;=\;\max\{\text{value of }M:\ M\text{ a matching of size }k\},
  \qquad 0\le k\le\lfloor n/2\rfloor .
\]
We usually drop the subscript. Note that $f(0)=0$.
\end{definitionbox}

\subsection{Which vertex sets can be covered?}
\begin{theorembox}[Covered sets]
A set $S\subseteq\{0,\dots,n-1\}$ is the set of covered vertices of some
matching if and only if every maximal run of consecutive elements of $S$
has even length. The matching is then unique.
\end{theorembox}
\begin{proof}
Let $M$ cover $S$ and take a maximal run $r,r+1,\dots,r+\ell-1$ of $S$.
Vertex $r-1\notin S$ is uncovered, so $r$ cannot be matched to the left;
its only other neighbour is $r+1$, so $\{r,r+1\}\in M$. Removing these two
vertices and repeating shows that the run is split into the pairs
$\{r,r+1\},\{r+2,r+3\},\dots$ --- hence $\ell$ is even, and $M$ is forced.
Conversely, if all runs are even, these pairs form a matching covering
exactly $S$.
\end{proof}

So $f(k)=\max\{x(S):|S|=2k,\ \text{all runs of }S\text{ even}\}$, where
$x(S)=\sum_{i\in S}x_i$. In particular, the \emph{uncovered} vertices
of a matching, read from left to right, split the covered ones into even
blocks; this is the picture used all the time below.

\subsection{Edge weights versus vertex weights}
Many problems put the weights on \emph{edges}: a path with vertices
$0,\dots,n$ and edge weights $w_0,\dots,w_{n-1}$, where $w_i$ belongs to
$\{i,i+1\}$, and one maximises the total weight of the chosen edges.

\begin{theorembox}[Edge weights $\leftrightarrow$ vertex weights]
Given $w_0,\dots,w_{n-1}$ put $x_0=0$ and $x_{i+1}=w_i-x_i$. Then
$w_i=x_i+x_{i+1}$ for all $i$, so every matching has the same total
edge weight and covered-vertex weight. Conversely, vertex weights $x$
give edge weights $w_i=x_i+x_{i+1}$ with the same property.
\end{theorembox}
\begin{proof}
The edge $\{i,i+1\}$ covers exactly the vertices $i$ and $i+1$, and
distinct edges of a matching cover distinct vertices. Hence
$\sum_{e\in M}w_e=\sum_{i\text{ covered}}x_i$ as soon as
$w_i=x_i+x_{i+1}$. The recursion is just this equation solved for
$x_{i+1}$.
\end{proof}
(The $x_i$ are alternating sums of the $w$'s, so they stay within
$n\cdot\max|w_i|$ in absolute value; 64-bit integers are enough in
practice.) From now on we work with vertex weights only.

\section{Three Problems, One Structure}\label{sec:reductions}
\subsection{Non-adjacent picks}
Let $B_0,\dots,B_{N-1}$ be the array. Build the path with $N+1$ vertices and
edge weights $w_i=B_i$. Two edges of this path share a vertex exactly
when their indices are neighbours, so a set of pairwise non-adjacent
entries is the same thing as a matching, with the same total weight.
Hence the answer for $K$ picks ($0\le K\le\lceil N/2\rceil$) is $f_x(K)$ with
$x$ obtained from $w=B$ by the previous theorem.

\subsection{Trades and subarrays}
The two other problems are one problem in two guises. Given
$A_0,\dots,A_{N-1}$ define prefix sums $P_0=0,\ P_i=A_0+\dots+A_{i-1}$,
and read them as prices. The subarray $[l,r)$ has sum $P_r-P_l$:
exactly the profit of buying at the end of day $l$ and selling on the
morning of day $r$. Subarrays $[l_1,r_1),\dots,[l_K,r_K)$ are disjoint
exactly when $l_1<r_1\le l_2<r_2\le\dots\le l_K<r_K$, which is exactly
the rule ``sell, then possibly buy again on the same day''. Conversely
prices give the array $A_i=P_{i+1}-P_i$. So it suffices to solve:

\begin{quote}
\textbf{$K$-trade problem.} Given $P_0,\dots,P_N$, maximise
$\sum_{j=1}^K(P_{r_j}-P_{l_j})$ over $0\le l_1<r_1\le l_2<r_2\le\dots\le l_K<r_K\le N$.
\end{quote}

Create a path with $2N$ vertices
\[
  v_0,v_1,\dots,v_{2N-1}\;=\;B_0,\,S_1,\,B_1,\,S_2,\,B_2,\,\dots,\,S_{N-1},\,B_{N-1},\,S_N,
\]
i.e.\ $v_{2l}=B_l$ is ``buy on day $l$'' and $v_{2r-1}=S_r$ is ``sell on
day $r$''; give $B_l$ the weight $+P_l$ and $S_r$ the weight $-P_r$.

\begin{figure}[H]
\centering
\begin{tikzpicture}[x=1.7cm,y=1cm]
\draw[thick,lectureink!70] (0,0)--(5,0);
\draw[line width=2.4pt,lectureaccent] (3,0)--(4,0);
\foreach \i/\lab/\w/\st in {0/B_0/P_0/lhi2, 1/S_1/-P_1/lhi2, 2/B_1/P_1/lhi2,
                             3/S_2/-P_2/lhi, 4/B_2/P_2/lhi, 5/S_3/-P_3/lhi2}{
  \node[lvert,\st] at (\i,0) {$\lab$};
  \node[llabel] at (\i,-0.75) {$\w$};
}
\node[llabel,anchor=west] at (-0.4,1.0) {orange = uncovered = the trades (buy $B_0$, sell $S_1$, buy $B_1$, sell $S_3$)};
\end{tikzpicture}
\caption{The path for $N=3$ and a matching of size $N-K=1$. The $2K=4$
uncovered vertices alternate buy, sell, buy, sell, and read as the two trades
$0\to1$ and $1\to3$.}\label{fig:trade}
\end{figure}

\begin{theorembox}[Trades $=$ matchings of size $N-K$]
A set $U\subseteq\{v_0,\dots,v_{2N-1}\}$ of $2K$ vertices is the uncovered
set of a matching of size $N-K$ if and only if, listed in path order,
it consists of a buy, a sell, a buy, a sell, \dots, i.e.\ it is
$B_{l_1},S_{r_1},B_{l_2},S_{r_2},\dots,B_{l_K},S_{r_K}$ with
$l_1<r_1\le l_2<r_2\le\dots$ --- a valid $K$-trade schedule. Moreover
\[
  (\text{weight covered by the matching})=P_0-P_N+(\text{profit of the schedule}),
\]
so the best profit with exactly $K$ trades is $f(N-K)-f(N)$.
\end{theorembox}
\begin{proof}
By the Covered-sets theorem, $U$ is an uncovered set iff the gaps before the
first element of $U$, between consecutive elements, and after the last one
all have even length. In terms of positions $p_1<\dots<p_{2K}$ in the
path $0,\dots,2N-1$ this says: $p_1$ is even, $p_{t+1}-p_t$ is odd, and
$2N-1-p_{2K}$ is even. So the positions alternate even, odd, even,
\dots, ending odd. Buys sit at even positions and sells at odd ones, hence
this is exactly ``buy, sell, \dots, buy, sell''. For the day numbers:
$2l<2r-1\iff l<r$ and $2r-1<2l'\iff r\le l'$, which are the schedule
rules.

For the weights, a schedule contributes $\sum_j(P_{l_j}-P_{r_j})=-\text{profit}$
to the weight of $U$. The whole path weighs $\sum_{i<N}P_i-\sum_{i\ge1}P_i=P_0-P_N$.
The covered weight is the total minus the weight of $U$, which gives the
formula. Finally $f(N)=P_0-P_N$ is the weight of the unique perfect
matching, so $f(N-K)-f(N)$ is the best profit.
\end{proof}

\section{The Greedy Algorithm}\label{sec:greedy}
\subsection{Augmenting one step at a time}
For an arbitrary graph with edge weights, let $M$ be a matching. An
\emph{$M$-augmenting path} is a path whose two end vertices are uncovered
by $M$ and whose edges alternately lie outside $M$ and in $M$, starting
and ending outside. Flipping it (taking the symmetric difference
$M\triangle C$) gives a matching with one more edge.

\begin{theorembox}[Successive optimal matchings]
Let $M$ be a maximum-weight matching among all matchings of size $k$, and
suppose a matching of size $k+1$ exists. Then there is an $M$-augmenting
path $C$ such that $M\triangle C$ has maximum weight among matchings of
size $k+1$. Consequently, if $C^*$ is an $M$-augmenting path maximising
$w(M\triangle C)$, then $M\triangle C^*$ is optimal for size $k+1$.
\end{theorembox}
\begin{proof}
Let $N$ be an optimal matching of size $k+1$. In $M\triangle N$ every
vertex has at most one edge from each matching, so the components are
paths and cycles whose edges alternate between $M$ and $N$. For a
component $C$ let $d(C)=|N\cap C|-|M\cap C|$; it is $0$ for cycles
and in $\{-1,0,1\}$ for paths. As $\sum_Cd(C)=|N|-|M|=1$, some path
component $C$ has $d(C)=+1$, so it starts and ends with $N$-edges.

\emph{$C$ is augmenting.} Let $z$ be an end vertex of $C$. It has no
$M$-edge on $C$. An $M$-edge at $z$ outside $C$ cannot lie in $M\triangle N$
(it would belong to the component $C$), so it would lie in $M\cap N$ and
then $z$ would have two different $N$-edges. Hence $z$ is uncovered by
$M$.

\emph{$M\triangle C$ is optimal.} A vertex of $C$ cannot lie on an edge of
$M\cap N$ (the same argument), so all $M$- and $N$-edges at vertices of
$C$ belong to $C$, and therefore $M\triangle C$ and $N\triangle C$ are
matchings, of sizes $k+1$ and $k$. Optimality of $M$ gives
$w(N\triangle C)\le w(M)$ and optimality of $N$ gives
$w(M\triangle C)\le w(N)$. But swapping the edges of $C$ between the two
matchings does not change the sum:
$w(M\triangle C)+w(N\triangle C)=w(M)+w(N)$. Both inequalities are
therefore equalities, and $M\triangle C$ is optimal.

For the last sentence: every augmenting path $C'$ yields a matching of
size $k+1$, hence $w(M\triangle C')\le f(k+1)$, while some $C$
attains $f(k+1)$; so the maximiser $C^*$ does too.
\end{proof}

\subsection{What augmenting paths look like on a path}
\begin{theorembox}[Augmenting paths on a path]
Let $M$ be a matching of the path. The $M$-augmenting paths are exactly
the segments $u,u+1,\dots,v$ where $u<v$ are \emph{consecutive uncovered
vertices} (no uncovered vertex strictly between them). With vertex weights
the flip changes the value by exactly $x_u+x_v$.
\end{theorembox}
\begin{proof}
An alternating path has an $M$-edge at each interior vertex, so interior
vertices are covered, so no uncovered vertex lies between the ends.
Conversely let $u<v$ be consecutive uncovered. The vertices strictly
between are covered and form a maximal run, which has even length by the
Covered-sets theorem and is matched as $\{u+1,u+2\},\{u+3,u+4\},\dots$.
So the segment alternates: non-$M$, $M$, non-$M$, \dots, non-$M$. After the flip the
matching is $\{u,u+1\},\{u+2,u+3\},\dots,\{v-1,v\}$, whose covered set is
the old one plus $\{u,v\}$. Only $u$ and $v$ are newly covered.
\end{proof}

\begin{figure}[H]
\centering
\begin{tikzpicture}[x=1.15cm,y=1cm]
\foreach \y in {0,1.9}{ \draw[thick,lectureink!35] (0,\y)--(7,\y); }
% before (y=1.9): M = {1,2},{4,5}
\draw[line width=2.4pt,lectureaccent] (1,1.9)--(2,1.9) (4,1.9)--(5,1.9);
\foreach \i/\st in {0/{},1/lhi,2/lhi,3/{},4/lhi,5/lhi,6/{},7/{}}{
  \node[lvert,\st] at (\i,1.9) {\i};
}
% after (y=0): M = {1,2},{3,4},{5,6}
\draw[line width=2.4pt,lectureaccent] (1,0)--(2,0) (3,0)--(4,0) (5,0)--(6,0);
\foreach \i/\st in {0/{},1/lhi,2/lhi,3/lhi2,4/lhi,5/lhi,6/lhi2,7/{}}{
  \node[lvert,\st] at (\i,0) {\i};
}
\draw[decorate,decoration={brace,amplitude=5pt}] (2.75,2.5)--(6.25,2.5)
  node[midway,above=5pt,llabel] {augmenting path $3,4,5,6$};
\draw[larrow] (4.5,1.35)--(4.5,0.6);
\node[llabel,anchor=east] at (-0.6,1.9) {before};
\node[llabel,anchor=east] at (-0.6,0) {after};
\end{tikzpicture}
\caption{Uncovered vertices $3$ and $6$ are consecutive. Flipping the path
between them re-pairs the covered block $\{4,5\}$ and adds exactly the
two end vertices to the covered set.}\label{fig:augment}
\end{figure}

\subsection{The algorithm}
\begin{theorembox}[Greedy algorithm]
Keep the list of uncovered vertices in path order, initially all
vertices. Repeat $\lfloor n/2\rfloor$ times: among the \emph{adjacent pairs
of the list} choose one $(u,v)$ with largest $x_u+x_v$, record the
\emph{gain} $g_t=x_u+x_v$ and delete $u$ and $v$ from the list.
Then $f(k)=g_0+g_1+\dots+g_{k-1}$ for every $k$.
\end{theorembox}
\begin{proof}
Induction on $k$; $f(0)=0$ is the empty matching. If the matching $M_k$
built so far has value $f(k)$, the list is the sequence of uncovered
vertices of $M_k$, adjacent list entries are consecutive uncovered
vertices, and by the previous theorem the best augmenting path has gain
$\max(x_u+x_v)$ over them. The Successive-optimal-matchings theorem then
shows the flipped matching is optimal of size $k+1$.
\end{proof}

\paragraph{Implementation.} A doubly linked list stores the uncovered
vertices; a max-heap stores the candidate pairs $(x_u+x_v,u,v)$. After
deleting $u,v$ with list neighbours $a$ and $d$ push the new candidate
$(a,d)$. A popped entry is stale unless both vertices are still in the list
and still neighbours. Each of the $O(n)$ pushes and pops costs $O(\log n)$.
This is \texttt{greedy\_pairs} in the listings, and it yields every
$f(k)$ in $O(n\log n)$.

\paragraph{The same trick for edge weights.} On edge weights
$w_i=x_i+x_{i+1}$ the new candidate $(a,d)$ has the value
$a+d=(a+b)+(c+d)-(b+c)=w_{\text{left}}+w_{\text{right}}-w_{\text{chosen}}$:
\emph{take the largest edge weight, then replace it and its two
neighbours by one edge of weight ``left $+$ right $-$ chosen''}. This is the
familiar ``regret'' trick for picking non-adjacent elements; here it is
just a vertex-weighted deletion in disguise.

\begin{theorembox}[Gains never increase]
In the greedy algorithm $g_0\ge g_1\ge\dots\ge g_{\lfloor n/2\rfloor-1}$.
Hence $f$ is concave, and the best value over all matchings of size
\emph{at most} $k$ is $g_0+\dots+g_{k'-1}$ with $k'=\min(k,\#\{t:g_t>0\})$.
\end{theorembox}
\begin{proof}
Let $(b,c)$ be the chosen pair, with list neighbours $a$ (left) and $d$ (right)
when they exist. Maximality gives $b+c\ge a+b$ and $b+c\ge c+d$, i.e.\ $c\ge a$
and $b\ge d$. The only new candidate is $(a,d)$, with
$a+d\le c+d\le b+c$. All other candidates were present before and are
$\le b+c$. So the next maximum is at most $g_t$.
\end{proof}

\section{Local Maxima and Linear Time}\label{sec:linear}
The heap always takes the \emph{global} maximum. This section shows that
any \emph{locally} largest pair may be taken instead, in any order. That
freedom allows a one-pass stack algorithm.

Fix a list $S$ of weights. A \emph{candidate} is a pair of neighbouring
entries; its value is their sum. Two candidates are \emph{neighbours} if
they share an entry. A candidate is a \emph{local maximum} if its value is
at least the value of each of its (at most two) neighbouring candidates.
The global maximum is a local maximum. Write $G(S)$ for the multiset of
gains $g_0,\dots,g_{\lfloor|S|/2\rfloor-1}$ of the greedy algorithm on $S$;
by the correctness of the greedy algorithm these are the increments $f_S(k+1)-f_S(k)$, so
$G(S)$ does not depend on how ties are broken.

\begin{figure}[H]
\centering
\begin{tikzpicture}[x=2.3cm,y=1cm]
\foreach \i/\n in {0/a,1/b,2/c,3/d}{ \node[lvert] (t\i) at (\i,1.9) {$\n$}; }
\draw[thick,lectureink!70] (t0)--(t1) node[midway,above,llabel]{$a+b$};
\draw[line width=2.4pt,lectureaccent] (t1)--(t2) node[midway,above,llabel]{$b+c$};
\draw[thick,lectureink!70] (t2)--(t3) node[midway,above,llabel]{$c+d$};
\node[llabel,text=lectureaccent] at (1.5,1.25) {local maximum};
\node[lvert] (u0) at (0,0) {$a$}; \node[lvert] (u3) at (3,0) {$d$};
\draw[thick,lecturewarn] (u0)--(u3) node[midway,above,llabel]{$a+d\le\min(a+b,\,c+d)$};
\draw[larrow] (1.5,1.0)--(1.5,0.5);
\node[llabel,anchor=west] at (1.62,0.75) {delete $b,c$};
\end{tikzpicture}
\caption{Deleting a local maximum $(b,c)$ replaces its two neighbouring
candidates by one new candidate $(a,d)$ that is not larger than either.}\label{fig:peel}
\end{figure}

\begin{theorembox}[Deleting a local maximum]
Let $(b,c)$ be a local maximum of $S$ with neighbours $a$ and $d$ (either may
be missing), and let $S'=S\setminus\{b,c\}$. The candidates of $S'$ are those
of $S$ not touching $b,c$, plus $(a,d)$ if both exist, and
\[ a+d\;\le\;\min(a+b,\;c+d)\;\le\;b+c . \]
If $(b,c)$ is strictly larger than both neighbours, the first inequality is strict.
\end{theorembox}
\begin{proof}
$b+c\ge a+b$ and $b+c\ge c+d$ give $c\ge a$ and $b\ge d$. Hence
$a+d\le c+d$ (from $a\le c$) and $a+d\le a+b$ (from $d\le b$), and
$\min(a+b,c+d)\le b+c$ by assumption. The strict case is the
same with strict inequalities.
\end{proof}

\begin{theorembox}[Peeling]
For every local maximum $(b,c)$ of a list $S$,
\[ G(S)\;=\;\{b+c\}\ \uplus\ G(S\setminus\{b,c\}) . \]
Consequently the following process records exactly $G(S)$, whichever local
maximum is chosen at each step: while at least two entries remain, pick
\emph{any} local maximum, record its value and delete it.
\end{theorembox}
\begin{proof}
\emph{Generic case.} Assume all subsets of the vertices have different
weight sums. Then all candidates of every sublist have different values, so
``local maximum'' means strictly larger than both neighbours. Call the
process in the statement a \emph{run}; the greedy algorithm is a run, since the
global maximum is a local maximum.

(A) \emph{If $p\ne q$ are local maxima of $T$ then they are disjoint, and $q$
is a local maximum of $T\setminus p$.} Overlapping candidates are neighbours,
and neighbours cannot each be strictly larger than the other. Deleting
$p=(b,c)$ changes the neighbours of $q$ only if $q=(d,e)$ or $q=(y,a)$; then
the neighbour $(c,d)$ resp.\ $(a,b)$ of $q$ is replaced by $(a,d)$, which is
smaller by the previous lemma. So $q$ still beats all its neighbours.

(B) \emph{All runs on $T$ record the same multiset $R(T)$.} Induction on
$|T|$; the cases $|T|\le1$ are empty. Let one run start with $p$ and another
with $q\ne p$. By (A) and the induction hypothesis, the first records
$\{p\}\uplus R(T\setminus p)$, and since a run on $T\setminus p$ may
start with $q$ (by (A)), this is $\{p,q\}\uplus R(T\setminus\{p,q\})$. By
symmetry the second run records the same.

Now $G(S)=R(S)$ because greedy is a run. Choosing the run that begins with
$(b,c)$ gives $R(S)=\{b+c\}\uplus R(S')=\{b+c\}\uplus G(S')$.

\emph{General case (ties).} For $\varepsilon>0$ raise $x_b$ and $x_c$ by
$\varepsilon$ and then add a further tiny perturbation $<\varepsilon/8$ to
every weight so that the weights are generic (non-generic vectors lie on
finitely many hyperplanes). The pair $(b,c)$ is then a strict local maximum of
the perturbed list, so the generic case applies. Let $\varepsilon\to0$. Each
$f_S(k)$ is a maximum of finitely many linear functions of the weights,
hence continuous, so the sorted gain sequence is continuous in the weights;
the same holds for $S'$. The identity passes to the limit.
\end{proof}

\subsection{A stack does it in one pass}\label{sec:stack}
Process the weights from left to right, keeping a stack $s_1,\dots,s_m$ of
the surviving entries so far (together with the unread entries it is the
current list). We maintain the invariant

\begin{quote}
\textbf{(I)} none of the candidates $(s_j,s_{j+1})$ with $j\le m-2$ is a
local maximum of the list.
\end{quote}
(Their neighbours $(s_{j-1},s_j)$ and $(s_{j+1},s_{j+2})$ lie inside the
stack, so this does not depend on the unread entries.)

\begin{theorembox}[Stack algorithm]
\begin{enumerate}[leftmargin=1.8em,itemsep=1pt]
  \item Push the next entry $v$. Let the top of the stack be $\dots,a,b,c,v$.
        While the candidate $(b,c)$ is a local maximum --- that is,
        $c\ge a$ (if $a$ exists) and $b\ge v$ --- record $b+c$ and delete $b,c$;
        then look at the new top again.
  \item After the last entry, while at least two entries remain, record the
        sum of the top two entries and delete them.
\end{enumerate}
The recorded values are exactly $G(S)$, and the running time is $O(n)$.
\end{theorembox}
\begin{proof}
\emph{Invariant.} After pushing $v$, only the candidate $(s_{m-1},s_m)$ of
the old stack has acquired a right neighbour $(s_m,v)$; for all
other $j$ the neighbours were already in the stack. So testing it restores
(I) unless it is a local maximum. If it is, delete it; the new candidate is
$(s_{m-2},v)$ at the top, and the only candidate whose neighbours changed is
$(s_{m-3},s_{m-2})$, namely the one now below the top --- exactly the one tested next.
Hence the loop re-establishes (I).

\emph{Correctness.} Every deleted candidate is a local maximum of the
current list, so by the Peeling theorem the recorded multiset plus the gains of
the remaining list is always $G(S)$. For the final phase: when all entries
are read the list is the stack. Some candidate has the globally maximal
value; it is a local maximum, so by (I) it is the top pair. Deleting
it keeps (I) true, because only the topmost candidate is then
exempt. Repeating until fewer than two entries remain records all
gains.

\emph{Time.} Each entry is pushed once and deleted at most once, and each
test costs $O(1)$.
\end{proof}

\subsection{Consequences}
\begin{itemize}[leftmargin=1.5em,itemsep=2pt]
  \item \textbf{All $f(k)$.} Sort the $\lfloor n/2\rfloor$ recorded values in
        decreasing order and take prefix sums. By the gains-never-increase and peeling theorems the result equals the heap's output; the
        constant factor is smaller.
  \item \textbf{One $f(K)$ in $O(n)$.} $f(K)$ is the sum of the $K$ largest
        recorded values, and a selection algorithm (quickselect, or
        \texttt{select\_nth\_unstable} in Rust) finds them without sorting.
        So $K$ non-adjacent picks cost $O(N)$ in total.
  \item \textbf{$K$ trades in $O(N)$.} Since $f(N)$ is the sum of \emph{all}
        gains, the Trades theorem of Section~\ref{sec:reductions} gives
        \[ \text{profit}(K)=f(N-K)-f(N)=-\bigl(\text{sum of the }K\text{ smallest gains}\bigr), \]
        found by one selection. In words: the best profits come from the
        $K$ \emph{most negative} gains of the trade path.
  \item \textbf{``At most $K$'' versions.} By the gains-never-increase theorem drop
        the non-positive terms (for trades: the non-negative gains).
\end{itemize}

\section{Static Range Queries}\label{sec:range}
\begin{definitionbox}[Range queries]
Given $x_0,\dots,x_{n-1}$, answer $q$ queries $(L,R,K)$: using only the
vertices $L,\dots,R-1$ (the induced path), what is the maximum value of a
matching of size $K$, where $0\le K\le\lfloor(R-L)/2\rfloor$?
\end{definitionbox}
The queries $K=0$ (answer $0$) and $2K=R-L$ (answer $x_L+\dots+x_{R-1}$)
are trivial, so assume $1\le K$ and $2K<R-L$. Write $I=[L,R)$.

The idea is to run the greedy algorithm once on the \emph{whole} path, stop it
at a well-chosen moment, and observe that almost everything it has
covered inside $I$ may be kept; then finish with $O(1)$ greedy steps.

\subsection{Which covered vertices can be kept?}
Let $M$ be a maximum-weight matching of the whole path among those of its
size, with covered set $D$, and put $s=|D\cap I|$. Let $c$ be the number of
edges of $M$ with exactly one endpoint in $I$ (\emph{crossing} edges; $c\le2$,
at most one per boundary). If the left boundary is crossed, let
$u$ be the leftmost uncovered vertex of $I$ and
\[ B_L=\{L,\,L+2,\,\dots,\,u-1\}, \]
and let $b_L$ be a vertex of $B_L$ of smallest weight. (Parity: $L$ is matched to
$L-1$, the vertices $L+1,\dots,u-1$ are paired up consecutively because
$u$ is uncovered, so $u-L$ is odd and $u-1\in B_L$.) For the right
boundary let $v$ be the rightmost uncovered vertex of $I$,
$B_R=\{R-1,R-3,\dots,v+1\}$ and $b_R$ a lightest vertex of $B_R$.
We call $b_L,b_R$ the \emph{release vertices}.

\begin{figure}[H]
\centering
\begin{tikzpicture}[x=0.92cm,y=1cm]
\tikzset{covd/.style={lvert,lhi},outs/.style={lvert,lgray},rel/.style={lvert,draw=lecturewarn,line width=1.4pt},
         newv/.style={lvert,lhi2},freev/.style={lvert}}
\def\xs{{7,-2,9,1,-1,8,4,-5,-8,2,5}}
% boundaries
\draw[dashed,lectureink!60] (0.5,-0.7)--(0.5,4.2) (9.5,-0.7)--(9.5,4.2);
\node[llabel] at (0.5,4.45) {$L=1$}; \node[llabel] at (9.5,4.45) {$R=10$};
\foreach \y in {0,1.7,3.4}{ \draw[thick,lectureink!35] (0,\y)--(10,\y); }
% row 1: global state at time T
\draw[line width=2.4pt,lecturewarn] (0,3.4)--(1,3.4) (9,3.4)--(10,3.4);
\draw[line width=2.4pt,lectureaccent] (2,3.4)--(3,3.4) (5,3.4)--(6,3.4);
\foreach \i/\w/\st in {0/7/outs,1/-2/covd,2/9/covd,3/1/covd,4/-1/freev,5/8/covd,6/4/covd,7/-5/freev,8/-8/freev,9/2/covd,10/5/outs}{
  \node[\st] at (\i,3.4) {\w};
}
% row 2: after release
\draw[line width=2.4pt,lectureaccent] (2,1.7)--(3,1.7) (5,1.7)--(6,1.7);
\foreach \i/\w/\st in {0/7/outs,1/-2/rel,2/9/covd,3/1/covd,4/-1/freev,5/8/covd,6/4/covd,7/-5/freev,8/-8/freev,9/2/rel,10/5/outs}{
  \node[\st] at (\i,1.7) {\w};
}
% row 3: after greedy
\draw[line width=2.4pt,lectureaccent] (1,0)--(2,0) (3,0)--(4,0) (5,0)--(6,0) (8,0)--(9,0);
\foreach \i/\w/\st in {0/7/outs,1/-2/newv,2/9/covd,3/1/covd,4/-1/newv,5/8/covd,6/4/covd,7/-5/freev,8/-8/newv,9/2/newv,10/5/outs}{
  \node[\st] at (\i,0) {\w};
}
\node[llabel,anchor=east] at (-0.7,3.4) {state at time $T$};
\node[llabel,anchor=east] at (-0.7,1.7) {release $b_L,b_R$};
\node[llabel,anchor=east] at (-0.7,0) {2 greedy steps};
\end{tikzpicture}
\caption{The query $(L,R,K)=(1,10,4)$ on $x=(7,-2,9,1,-1,8,4,-5,-8,2,5)$.
Top: the global greedy has covered $\{0,1,2,3,5,6,9,10\}$, and the edges
$\{0,1\}$ and $\{9,10\}$ (orange) cross the boundaries. Middle: $u=4$, $v=8$,
$B_L=\{1,3\}$, $B_R=\{9\}$; the lightest are $b_L=1$ and $b_R=9$, which are
released. Bottom: the greedy on the uncovered vertices $1,4,7,8,9$ of $I$ picks $(1,4)$
and then $(8,9)$. Result $9+1+8+4-3-6=13$.}\label{fig:range}
\end{figure}

\begin{theorembox}[Release lemma]
Assume $I$ contains an uncovered vertex of $M$ and $s+c\le 2K$. Then some
maximum-weight matching of size $K$ inside $I$ covers every vertex of
$(D\cap I)\setminus\{b_L,b_R\}$. (A release vertex exists only on a side that
is crossed.)
\end{theorembox}
\begin{proof}
Let $M_I$ be the set of edges of $M$ that meet $I$; it has
$(s+c)/2\le K$ edges. Among all maximum-weight size-$K$ matchings $H$ inside
$I$ choose one with $|H\cap M_I|$ as large as possible. The graph
$M_I\triangle H$ is a union of alternating paths; put
$d(C)=|H\cap C|-|M_I\cap C|$ for a component $C$, and note
$\sum_Cd(C)=K-|M_I|\ge0$.

\emph{Components.} A component with a crossing edge is a \emph{boundary}
component; it starts at the outer endpoint with that $M$-edge (the outer
vertex has no $H$-edge), so $d\in\{-1,0\}$. No component contains both
crossing edges: it would contain the whole stretch $L-1,\dots,R$, in
particular $u$ as an interior vertex, but interior vertices of an alternating
path are covered by $M$. All other components are \emph{internal}, inside $I$.

\emph{Swapping.} Let $\mathcal A$ be a set of internal components with
$\sum_{\mathcal A}d=0$. Exchanging $M$ and $H$ on $\mathcal A$ gives two
matchings (internal vertices lie in $I$, so all their $M$-edges are in
$M_I$) that keep their sizes. As $M$ is optimal for its size and $H$ for its,
$w(H\cap\mathcal A)\le w(M_I\cap\mathcal A)\le w(H\cap\mathcal A)$, so the
two parts weigh the same and the swapped $H$ is again optimal --- and has more
edges in common with $M_I$ if $\mathcal A\ne\emptyset$ (such an $\mathcal A$
contains as many $M_I$-edges as $H$-edges, at least one). That contradicts
the choice of $H$. Hence no nonempty such $\mathcal A$ exists: all internal
components have the same $d\ne0$ (else take one of each sign). Since
boundary components have $d\le0$ and the total is $\ge0$, any internal
components all have $d=+1$; they begin and end with $H$-edges.

\emph{Which covered vertices does $H$ miss?} Let $z\in D\cap I$ be uncovered
by $H$. Its $M$-edge is not in $H$, so it lies in a component $C$ in which $z$
is an end vertex whose last edge is an $M$-edge. Internal components
($d=+1$) have no such end, so $C$ is a boundary component with $d=-1$.
There is at most one per side. Say $z$ is on the left side. $C$ is the path
$L-1,L,\dots,z$ with edges $M,H,M,\dots,M$; it cannot reach $u$, so
$z<u$, and the odd number of edges gives $z\equiv L\pmod2$; thus $z\in B_L$.
Inside $[L,u)$ the edges of $C$ give $H$-edges $\{L,L+1\},\dots,\{z-2,z-1\}$,
and the remaining vertices $z+1,\dots,u-1$ are paired by $M$ as
$\{z+1,z+2\},\dots,\{u-2,u-1\}$. These pairs belong to $H$ too: a pair
outside $H$ would lie in an internal component, an interval
with $H$-edges at both ends, whose end vertices are uncovered by $M$ --- but
all uncovered vertices of $I$ are $\ge u$, a contradiction.
(It cannot be in the right boundary component either, as that would contain
$u$ as an interior vertex.)

\emph{Moving the hole.} In $H$ replace the edges inside $[L,u)$ by
$\{L,L+1\},\dots,\{b-2,b-1\},\{b+1,b+2\},\dots,\{u-2,u-1\}$ with
$b=b_L$; this is valid because $b\equiv L\pmod2$, and it has the same number
of edges. It covers $[L,u)\setminus\{b\}$ instead of $[L,u)\setminus\{z\}$,
changing the weight by $x_z-x_b\ge0$ because $b_L$ is lightest in $B_L$
and $z\in B_L$. So the new matching is still optimal. Do the same on the right.
Every vertex of $D\cap I$ other than $b_L,b_R$ is now covered.
\end{proof}

\subsection{Reducing to a greedy on the rest}
After releasing, put $D'=(D\cap I)\setminus\{b_L,b_R\}$. Its maximal runs
inside $I$ are even: runs not touching a released vertex are runs of
$M$ inside $I$; on the left, $[L,u)\setminus\{b_L\}$ splits into
$[L,b_L)$ and $(b_L,u)$ of lengths $b_L-L$ and $u-b_L-1$, both even; the
right side is symmetric. Let $F=I\setminus D'$ be the \emph{free}
vertices, in path order.

\begin{theorembox}[Finish by greedy]
The maximum value of a size-$K$ matching inside $I$ is
\[ x(D')\;+\;f_F\bigl(K-|D'|/2\bigr), \]
where $f_F$ is the function for the weight sequence $F$ (a path on the free
vertices, adjacent in $F$ even if not in the original path).
\end{theorembox}
\begin{proof}
By the Release lemma the optimum is attained by a matching whose covered set
is $Q\supseteq D'$; write $T=Q\setminus D'\subseteq F$. A maximal run of $Q$
consists of runs of $D'$ (even length) and elements of $T$; the corresponding
run of $T$ in the list $F$ has the same parity. So $Q$ has only even runs iff
$T$ has only even runs in $F$, and the optimum is $x(D')$ plus the best
$T$ with $|T|=2K-|D'|$, i.e.\ $f_F(K-|D'|/2)$.
\end{proof}
By the greedy theorem $f_F(m)$ is the sum of the first $m$ gains of
the greedy algorithm run on $F$. Hence the whole query is: release at most two
vertices, then perform $m=K-|D'|/2$ greedy steps on the free vertices of
$I$, and the answer is the total weight of the covered vertices inside $I$
afterwards. It remains to show that $m\le2$ and that the hypotheses of the
lemma hold, for a suitable state of the global greedy.

\subsection{The focus time}
Run the greedy algorithm of Section~\ref{sec:greedy} on the whole path; step
$t=0,1,\dots$ covers two vertices, and we say these vertices have
\emph{time} $t$ (if $n$ is odd, the one never covered gets time $\lfloor n/2\rfloor$).
Let $D_t$ be the set of vertices covered after $t$ steps (those with time
$<t$). Its unique matching $M_t$ has value $f(t)$, so it is
optimal for its size, as the Release lemma needs.

For a query let the \emph{focus time} $T$ be the $2K$-th smallest time
among the vertices of $I$. Then $D_T$ contains fewer than $2K$ vertices
of $I$ while $D_{T+1}$ contains at least $2K$; as a step adds at most two
vertices of $I$,
\[ s=|D_T\cap I|\in\{2K-2,\,2K-1\}. \]
The crossing edges are determined by parities. At most one run of $D_T$
contains both $L-1$ and $L$; all other runs lie on one side of the boundary
and have even length. If the straddling run has $p$ vertices left of $L$, it is
paired as $(1,2),(3,4),\dots$ and the pair $\{L-1,L\}$ occurs iff $p$ is odd.
So the left boundary is crossed iff $|D_T\cap[0,L)|$ is odd, and likewise on
the right with $|D_T\cap[R,n)|$. Since $s\equiv c\pmod 2$ (the total $|D_T|$ is even) and $c\le2$,
\begin{center}\small\sffamily
\begin{tabular}{@{}cccc@{}}
\toprule
$s$ & $c$ & released vertices & greedy steps $m=K-|D'|/2$\\
\midrule
$2K-2$ & $0$ & $0$ & $1$\\
$2K-1$ & $1$ & $1$ & $1$\\
$2K-2$ & $2$ & $2$ & $2$\\
\bottomrule
\end{tabular}
\end{center}
In every case $s+c\le2K$, and $s\le2K-1<R-L$ guarantees an uncovered vertex in
$I$. So the Release lemma applies, and each query needs one or two greedy
steps.

\subsection{The offline algorithm and its data structures}
Sort the queries by focus time and sweep $t=0,1,\dots$ over the global
greedy; at time $t$ the structures represent $D_t$. For each query with
focus time $t$:
\begin{enumerate}[leftmargin=1.8em,itemsep=1pt]
  \item Count the covered vertices left of $L$ and right of $R-1$ to detect
        crossings; find $u$ and/or $v$; take the minimum-weight vertex of the
        right parity in $[L,u)$ and/or $(v,R)$ and \emph{uncover} it.
  \item Repeat $m$ times: find the adjacent uncovered pair inside $I$ with the
        largest sum and cover it.
  \item The answer is the total weight of the covered vertices inside $I$.
  \item Undo steps 1--2 in reverse order.
\end{enumerate}
Then apply global step $t$ and continue. The structures needed are:

\begin{itemize}[leftmargin=1.5em,itemsep=1pt]
  \item The focus time of $(L,R,K)$ is the $(2K-1)$-st value (0-indexed) of
        the array of times restricted to $[L,R)$: a \textbf{wavelet matrix}
        built once, $O(\log n)$ per query.
  \item A \textbf{Fenwick tree of 0/1 flags} (uncovered) gives counts of
        covered vertices in a range, the next/previous uncovered vertex by
        binary descent, and a second Fenwick tree of the weights of
        covered vertices gives the final sum.
  \item Two \textbf{sparse tables}, one for even and one for odd indices,
        give the lightest vertex of a given parity in a range.
  \item A \textbf{segment tree} indexed by $i$ storing $x_i+x_{\mathrm{next}(i)}$
        for each uncovered $i$ (neighbour in the list of uncovered vertices)
        gives the best adjacent uncovered pair inside $I$: query the range
        $[L,p)$, where $p$ is the last uncovered vertex of $I$, so that
        the right element of the pair is also in $I$.
\end{itemize}
Covering or uncovering a vertex changes $O(1)$ entries of these
structures, each in $O(\log n)$. Each query makes $O(1)$ changes and
undoes them, so the total time is $O((n+q)\log n)$ after the
$O(n\log n)$ global greedy.

\begin{notebox}[On ties]
Neither the heap, the focus-time argument nor the Release lemma needs
distinct weights: the Release lemma is proved for arbitrary weights (the
extremal choice of $H$ replaces genericity), and the greedy theorem allows
any maximiser at each step. The reference code uses whatever tie-breaking
its heap, sparse tables and segment tree happen to give, and the tests
below deliberately use weights in $\{-1,0,1\}$ to force many ties.
\end{notebox}

\section{Reference Implementations}\label{sec:code}
Both listings contain the heap greedy, the stack algorithm with selection,
the reductions of Section~\ref{sec:reductions} (non-adjacent picks, $K$
trades, $K$ disjoint subarrays) and the offline range-query structure.
They are cross-checked against dynamic programming on thousands of random
inputs, including the tie-heavy small-range ones: all three ways of computing
$f$ agree, the reductions match direct DPs for non-adjacent picks and for
exactly $K$ subarrays, and about $10^5$ range queries per language
agree with a DP on the slice. The harness also confirmed that the
cases $(s-2K,c)\in\{(-2,0),(-2,2),(-1,1)\}$ of the table above are the
only ones that occur.

\codenote
\subsection{Rust}
\corelisting{Rust}{example.rs}

\subsection{Python}
\corelisting{Python}{example.py}

\section{Worked Examples}
\begin{examplebox}[Stack algorithm]
Take $x=(3,-1,4,1,-5,9,2,-6)$. Pushing $3,-1,4,1$ changes nothing: with
$a,b,c,v=3,-1,4,1$ the candidate $(-1,4)$ has $c=4\ge a=3$ but $b=-1<v=1$. Pushing
$-5$ gives the top $-1,4,1,-5$: here $(4,1)$ satisfies $1\ge-1$ and
$4\ge-5$, so record $5$ and delete; the stack is $3,-1,-5$, where $(3,-1)$ is
at the list start and $3\ge-5$: record $2$, stack $-5$. Push $9,2$ and
then $-6$: the top is $-5,9,2,-6$ and $(9,2)$ is a local maximum:
record $11$. The stack is $-5,-6$; the final phase records $-11$. The
multiset of gains is $\{11,5,2,-11\}$, so
\[ f(0..4)=0,\ 11,\ 16,\ 18,\ 7 . \]
The heap greedy takes the same pairs in the order $(9,2),(4,1),(3,-1),(-5,-6)$.
\end{examplebox}

\begin{examplebox}[Non-adjacent picks]
$B=(5,-2,6,4,-1,3)$ becomes $x=(0,5,-7,13,-9,8,-5)$ and the greedy gains are
$6,5,3,\dots$, so the best sums for $K=1,2,3$ are $6,\ 11,\ 14$
(for $K=3$: indices $0,2,5$ give $5+6+3$).
\end{examplebox}

\begin{examplebox}[Stock trading]
Prices $P=(5,2,4,1,6,3,7)$, $N=6$. The trade path has weights
$5,-2,\,2,-4,\,4,-1,\,1,-6,\,6,-3,\,3,-7$ and the smallest gains of
the matching problem are $-6,-3,-2,\dots$. The best profits for
$K=1,2,3$ are therefore $6,\ 9,\ 11$: the single best trade is buy at $1$,
sell at $7$; three trades are $2\to4$, $1\to6$ and $3\to7$. For
$K=4$ the remaining gains are positive, so the profit drops to $8$: you
are \emph{forced} to make four trades.
\end{examplebox}

\begin{examplebox}[Range query]
Figure~\ref{fig:range} follows $x=(7,-2,9,1,-1,8,4,-5,-8,2,5)$ and
the query $(L,R,K)=(1,10,4)$. The global greedy picks
$(5,6),(2,3),(9,10),(0,1),\dots$; the times of vertices $1,\dots,9$ are
$3,1,1,4,0,0,4,\ast,2$, so the focus time is $T=4$ and $D_4=\{0,1,2,3,5,6,9,10\}$ gives
$s=6=2K-2$ and $c=2$. Release $b_L=1$ and $b_R=9$; the free vertices
are $1,4,7,8,9$; two greedy steps take $(1,4)$ with sum $-3$ and then
$(8,9)$ with sum $-6$. The answer is
$x(D')+(-3)+(-6)=9+1+8+4-9=13$, which a direct DP on $x_1..x_9$ confirms.
\end{examplebox}

\section{Practice Problems}
\begin{itemize}[leftmargin=1.5em,itemsep=1pt]
  \item JOISC 2018 Day 4 \emph{Candies} --- pick exactly $K$ non-adjacent
        elements for every $K$ up to $\lceil N/2\rceil$; the standard
        test for the heap greedy of Section~\ref{sec:greedy}.
  \item APIO/CTSC 2007 \emph{Data Backup} --- pair up $2K$ of the points on a
        line to minimise total cable length: non-adjacent gaps, negated.
  \item Codeforces 958E2 \emph{Guard Duty (medium)} --- the same idea with
        larger limits.
  \item Codeforces 391F3 \emph{Stock Trading} --- at most $K$ trades with
        $N$ up to $4\cdot10^6$; use the linear-time algorithm of
        Section~\ref{sec:linear}.
  \item Codeforces 280D \emph{k-Maximum Subsequence Sum} --- the dynamic
        cousin: updates and range queries for $\le K$ disjoint subarrays.
        The static version is Section~\ref{sec:range}.
  \item Luogu P1484 \emph{Planting Trees} --- non-adjacent picks, the
        simplest form.
\end{itemize}

\section{Further Reading}
\begin{itemize}[leftmargin=1.5em]
  \item AtCoder Algorithm Lectures index (source of this curriculum):
        \url{https://info.atcoder.jp/entry/algorithm_lectures/index}
  \item M.\ Cs\H{u}r\"os, ``Maximum-scoring segment sets'', \emph{IEEE/ACM
        Transactions on Computational Biology and Bioinformatics} 1(4),
        2004, pp.\ 139--150.
  \item F.\ Bengtsson and J.\ Chen, \emph{Computing maximum-scoring
        segments optimally}, Lule\aa{} University of Technology Research
        Report 2007:3, 2007 (a linear-time algorithm for $K$ disjoint
        maximum-sum segments).
  \item L.\ Lov\'asz and M.\ D.\ Plummer, \emph{Matching Theory}, for the
        general theory of augmenting paths in weighted matching.
\end{itemize}

\end{document}
